%!TEX encoding = IsoLatin

%
% Chapitre "Structure d'un rapport technique"
%

\chapter{Description}
\label{s:structure_rapport}
L'équipe NordiQc a reçu la tâche de faire le design conceptuel d’un système. Ce système doit être en mesure de suivre l’activité, de compter et de documenter la faune arctique dans le Grand Nord canadien et ce en minimisant son impact environnemental. Le système doit rester opérationnel pendant une durée minimale d’un an. Durant cette période, il doit être en mesure de prélever et d’enregistrer les données recueillies de manière autonome. L’appareil doit être capable de fonctionner dans des conditions hivernales et estivales, soit des températures de ±20° Celsius et doit rester opérationnel sous deux mètres de neige. L’appareil doit comprendre une caméra ayant au minimum une résolution de 720p à 15 images par seconde. Les vidéos doivent être d’un minimum de cinq secondes et le délai maximal entre chaque enregistrement est de 1h. \\

 L’appareil doit être facilement transportable et déployable par deux personnes qui n’ont pas de formation en ingénierie. Le système doit pouvoir communiquer avec la centrale à une distance de plus de 2000 km. De plus, celui-ci doit aussi renseigner sur l’état du système et le nombre d’événements enregistrés. L’utilisateur du système doit pouvoir accéder aux données de la centrale depuis une application mobile. Cette dernière doit être sécurisée pour trois personnes. Finalement, le coût et la masse de l’appareil sont aussi des enjeux à minimiser.  




